% !PRINTED 153
\weberpage{153}

% [restored 2026-08-06] omission raised by this page's own translation pass and confirmed
% against the scan; the German fragment now carries the line too.
Herein one is to set, according to \S.\,42,
\[
\begin{aligned}
s_0 &= 3, \quad a_0 s_1 = -a_1\\
a_0^2 s_2 &= a_1^2 - 2 a_0 a_2\\
a_0^3 s_3 &= -a_1^3 + 3 a_0 a_1 a_2 - 3 a_0^2 a_3\\
a_0^4 s_4 &= a_1^4 - 4 a_0 a_1^2 a_2 + 4 a_0^2 a_1 a_3 + 2 a_0^2 a_2^2,
\end{aligned}
\]
whereby one obtains

\weq{10}{D = a_1^2 a_2^2 + 18 a_0 a_1 a_2 a_3 - 4 a_0 a_2^3 - 4 a_1^3 a_3 - 27 a_0^2 a_3^2.}

\begin{center}
\S.\,47.

\gesperrt{Discriminants of the Forms of Third and Fourth Order.}
\end{center}

For the functions of third and fourth order we have, in the solutions of the equations of these two degrees (\S.\,35, 37), a simple means of computing the discriminants.

As regards first the equation of third order, the three roots of
\weq{1}{x^3 + a x + b = 0}
were represented in the form:
\weq{2}{\begin{aligned}
\alpha &= u + v\\
\beta &= \varepsilon u + \varepsilon^2 v\\
\gamma &= \varepsilon^2 u + \varepsilon v,
\end{aligned}}
in which
\weq{3}{\varepsilon = \frac{-1 + \sqrt{-3}}{2}, \quad \varepsilon^2 = \frac{-1 - \sqrt{-3}}{2}}
are the complex third roots of unity, and $u$, $v$ are determined by \S.\,35, (8).

From (3) it follows that
\weq{4}{\varepsilon - \varepsilon^2 = \sqrt{-3}, \quad 1 - \varepsilon = \varepsilon^2 \sqrt{-3}, \quad 1 - \varepsilon^2 = -\varepsilon \sqrt{-3}}
and from (2)
\[
\begin{aligned}
\alpha - \beta &= \sqrt{-3}\,(\varepsilon^2 u - \varepsilon v)\\
\alpha - \gamma &= -\sqrt{-3}\,(\varepsilon u - \varepsilon^2 v)\\
\beta - \gamma &= \sqrt{-3}\,(u - v).
\end{aligned}
\]
From this, by multiplication,
\weq{5}{(\alpha - \beta)(\alpha - \gamma)(\beta - \gamma) = 3 \sqrt{-3}\,(u^3 - v^3).}
